% Part: proof-theory
% Chapter: cut-elimination
% Section: ce-largest

\documentclass[../../../include/open-logic-section]{subfiles}

\begin{document}

\olfileid[id]{pt}{cut}{inv}

\olsection{Menghapus Potongan Berperingkat Terbesar}

Bukti \olref[seq][inv]{lem:G3c-invert} menunjukkan bahwa jika \Log{G3c}
!!{prove} kesimpulan suatu aturan logis (selain \RightR{\lexists} dan
\LeftR{\lforall}), maka sistem itu juga !!{prove} setiap premis aturan itu
tanpa menambah !!{height} bukti. Kita dapat menunjukkan hasil yang lebih
kuat: hal ini berlaku pula untuk $\Log{G3c} + \CutCS$.

Seperti sebelumnya, kita mendefinisikan \emph{peringkat potongan} suatu
inferensi \CutCS{} sebagai kedalaman !!{formula} potongannya.

\begin{lem}\ollabel{lem:inv-G3c-cut}
Jika $\Log{G3c} + \CutCS$ !!{prove} kesimpulan suatu aturan logis~$R$
selain \RightR{\lexists} atau \LeftR{\lforall} dengan menggunakan
!!a{proof}~$\pi$, maka sistem itu juga !!{prove} setiap premis dengan
!!{proof}~$\pi'$ (atau !!{proof}s~$\pi'$, $\pi''$, bergantung pada apakah
aturan tersebut mempunyai satu atau dua premis). Selain itu,
(a)~$\pheight{\pi'}\le\pheight{\pi}$ (dan
$\pheight{\pi''}\le\pheight{\pi}$) serta (b)~banyaknya inferensi \CutCS{}
maupun peringkat masing-masing sama dalam~$\pi'$ (dan $\pi''$) seperti
dalam~$\pi$.
\end{lem}

\begin{proof}
Dengan memeriksa bukti
\cref{pt:seq:inv:lem:G3c-invert,pt:seq:inv:lem:invert-quant}. Bukti tersebut
berjalan melalui induksi pada~$\pheight{\pi}$. Dalam kasus dasar, kita
mengganti suatu aksioma dengan aksioma lain. Jadi, kondisi (a) dan~(b)
terpenuhi secara sepele. Jika aturan terakhir dalam~$\pi$ adalah~$R$,
sub-!!{proof}s langsungnya adalah !!{proof}s~$\pi_i$ yang kita inginkan,
dan (a) maupun~(b) terpenuhi. Dalam semua kasus lain, kita menerapkan
hipotesis induksi pada sub-!!{proof}s langsung dari~$\pi$, yakni yang
menuju ke premis-premis inferensi terakhir~$R$, lalu menerapkan aturan~$R$
yang sama pada hasilnya. Kondisi (a) dan~(b) berlaku bagi sub-!!{proof}s
langsung bukti baru, sehingga (a) dan~(b) berlaku bagi seluruh bukti. Tinggal
memeriksa kasus ketika inferensi terakhir adalah~\CutCS. Sekali lagi, kita
hanya membahas satu contoh: andaikan $\pi$ merupakan !!a{proof} dari
kesimpulan $!B \land !C, \Gamma \Sequent \Delta$ milik
\LeftR{\land}, dan berakhir dengan~\CutCS:
\[
\AxiomC{}
\RightLabel{$\pi_1$}
\Deduce$!B \land !C, \Gamma \fCenter \Delta, !A$
\AxiomC{}
\RightLabel{$\pi_2$}
\Deduce$!A, !B \land !C, \Gamma \fCenter \Delta$
\RightLabel{\CutCS}
\BinaryInf$!B \land !C, \Gamma \fCenter \Delta$
\DisplayProof
\]
Hipotesis induksi berlaku pada $\pi_1$ dan~$\pi_2$ (yang juga merupakan
!!{proof}s atas instansiasi kesimpulan $\LeftR{\land}$) dan menghasilkan
!!{proof}s $\pi_1'$ dan~$\pi_2'$ masing-masing dari $!B, !C, \Gamma
\fCenter \Delta, !A$ dan $!A, !B, !C, \Gamma \fCenter \Delta$. Kita
menggabungkannya dengan \CutCS{} untuk memperoleh~$\pi'$:
\[
\AxiomC{}
\RightLabel{$\pi_1'$}
\Deduce$!B, !C, \Gamma \fCenter \Delta, !A$
\AxiomC{}
\RightLabel{$\pi_2'$}
\Deduce$!A, !B, !C, \Gamma \fCenter \Delta$
\RightLabel{\CutCS}
\BinaryInf$!B, !C, \Gamma \fCenter \Delta$
\DisplayProof
\]
Jelaslah bahwa (a) dan~(b) terpenuhi.
\end{proof}

Perhatikan bahwa cara ini tidak akan berhasil jika kita menggunakan \Cut{}
alih-alih \CutCS, sebab kesimpulan !!{proof} terakhir itu akan mempunyai dua
salinan $!B, !C, \Gamma$ dalam anteseden dan dua salinan $\Delta$ dalam
suksedennya. Kita harus mengacu pada dapat diterimanya kontraksi, tetapi
bukti \olref[seq][inv]{prop:G3c-cont-adm} menggunakan lema keterbalikan.
Dengan demikian, kita akan berargumen secara melingkar.

Kita dapat menggunakan \cref{pt:cut:inv:lem:inv-G3c-cut} untuk memberikan
bukti alternatif eliminasi potongan. Dalam bukti ini, kita tidak secara
berturut-turut menghapus kemunculan teratas inferensi \CutCS{}, melainkan
kemunculan inferensi \CutCS{} yang berperingkat maksimal. Artinya, kita
memulai dengan !!a{proof}~$\pi$ dalam $\Log{G3c} + \CutCS$, lalu menghapus
suatu potongan di mana saja di dalamnya (mungkin dengan menggantinya dengan
potongan berperingkat lebih rendah)---dan terus melakukannya hingga tidak ada
potongan yang tersisa. Satu-satunya syarat ialah bahwa di atas potongan yang
kita tangani tidak terdapat potongan berperingkat sama atau lebih tinggi.

\begin{lem}\ollabel{lem:max-cut-red-G3c}
Jika $\pi$ adalah !!a{proof} dalam $\Log{G3c}+\CutCS$ yang berakhir dengan
inferensi~\CutCS{} berperingkat~$n$, dan selebihnya hanya menggunakan
aturan-aturan~$\Log{G3c}$ serta inferensi \CutCS{} berperingkat $<n$, maka
terdapat bukti~$\pi'$ atas sekuen akhir yang sama dalam~$\Log{G3c} +
\CutCS$, dengan setiap inferensi \CutCS{} berperingkat~$<n$.
\end{lem}

\begin{proof}
!!{proof}~$\pi$ berakhir dengan:
\[
\AxiomC{}
\RightLabel{$\pi_1$}
\Deduce$\Gamma \fCenter \Delta, !A$
\AxiomC{}
\RightLabel{$\pi_2$}
\Deduce$!A, \Gamma \fCenter \Delta$
\RightLabel{\CutCS}
\BinaryInf$\Gamma \fCenter \Delta$
\DisplayProof
\]
Kita membedakan kasus-kasus menurut bentuk~$!A$.

Andaikan $!A$ atomik. Perhatikan bahwa syarat semua inferensi \CutCS{} dalam
$\pi_1$ dan $\pi_2$ harus mempunyai peringkat $< \depth{!A} = 0$
menyiratkan bahwa tidak mungkin ada inferensi \CutCS{} dalam~$\pi_1$
ataupun~$\pi_2$. Kita menunjukkan bahwa $\Gamma \Sequent \Delta$ mempunyai
bukti bebas-potongan melalui induksi pada $\pheight{\pi_2}$. Jika
$\pheight{\pi_2} = 0$, maka $!A, \Gamma \Sequent \Delta$ adalah suatu
aksioma. Jika $!A$ tidak prinsipal, suatu $!C$ atomik merupakan
!!a{element} dari~$\Gamma$ maupun $\Delta$, dan $\Gamma \Sequent \Delta$
sendiri merupakan aksioma. Jika $!A$ prinsipal, maka $\Delta = \Delta',
!A$. Dalam hal ini, $\pi_1$ berakhir dengan $\Gamma \Sequent \Delta', !A,
!A$. Kita memperoleh bukti bebas-potongan dari $\Gamma \Sequent \Delta',
!A$ dengan menerapkan \olref[seq][inv]{prop:G3c-cont-adm} (dapat diterimanya
kontraksi). Jika $\pheight{\pi_2} > 0$, bukti itu berakhir dengan suatu
aturan~$R$. Ambil salah satu premis aturan tersebut: premis itu berbentuk
$!A, \Gamma', \Pi \Sequent \Delta', \Lambda$, dengan $\Pi$ dan $\Lambda$
sebagai !!{formula}s aktif dari~$R$. Penerapan
\olref[seq][adm]{prop:weak-G3c-adm} pada~$\pi_1$ dan $\pi_2$ menghasilkan
!!{proof}s dari $\Gamma, \Gamma', \Pi \Sequent \Delta, \Delta', \Lambda,
!A$ dan $!A, \Gamma, \Gamma', \Pi \Sequent \Delta, \Delta', \Lambda$,
yang memenuhi hipotesis induksi. Dengan demikian, kita memperoleh !!a{proof}
dari $\Gamma, \Gamma', \Pi \Sequent \Delta, \Delta', \Lambda$ untuk
setiap premis~$R$. Penerapan $R$ menghasilkan $\Gamma, \Gamma \Sequent
\Delta, \Delta$, lalu \olref[seq][inv]{prop:G3c-cont-adm} memberikan bukti
bebas-potongan dari $\Gamma \Sequent \Delta$.

Jika $!A$ tidak atomik, kita membedakan kasus-kasus menurut !!{main
operator} dari~$!A$. Dalam setiap kasus, kita menerapkan
\olref{lem:inv-G3c-cut} untuk memperoleh !!{proof}s atas premis-premis aturan
kiri dan kanan dengan $!A$ sebagai formula prinsipal. Lalu kita melanjutkan
seperti dalam kasus~E pada bukti \olref[top]{lem:cut-adm-G3c}.
\begin{enumerate}
  \item $!A \ident \lnot !B$. !!{proof}s $\pi_1$ dan $\pi_2$ berakhir
  dengan $\Gamma \Sequent \Delta, \lnot !B$ dan $\lnot !B, \Gamma \Sequent
  \Delta$. Menurut \olref{lem:inv-G3c-cut}, terdapat !!{proof}s $\pi_1'$
  dan $\pi_2'$ dari $!B, \Gamma \Sequent \Delta$ dan $\Gamma \Sequent
  \Delta, !B$. !!{proof}~$\pi'$ ialah:
  \[
  \AxiomC{}
  \RightLabel{$\pi_2'$}
  \Deduce$\Gamma \fCenter \Delta, !B$
  \AxiomC{}
  \RightLabel{$\pi_1'$}
  \Deduce$!B, \Gamma \fCenter \Delta$
  \RightLabel{\CutCS}
  \BinaryInf$\Gamma \fCenter \Delta$
  \DisplayProof
  \]
  \item $!A \ident (!B \lor !C)$. !!{proof}s $\pi_1$ dan $\pi_2$
  berakhir dengan $\Gamma \Sequent \Delta, !B \lor !C$ dan $!B \lor !C,
  \Gamma \Sequent \Delta$. Menurut \olref{lem:inv-G3c-cut} (keterbalikan
  untuk \RightR{\lor} yang diterapkan pada~$\pi_1$), terdapat !!a{proof}
  $\pi_1'$ dari $\Gamma \Sequent \Delta, !B, !C$. Menurut
  \olref{lem:inv-G3c-cut} (keterbalikan untuk \LeftR{\lor} yang diterapkan
  pada~$\pi_2$), terdapat !!a{proof} $\pi_2'$ dari $!B, \Gamma \Sequent
  \Delta$ dan !!a{proof} $\pi_2''$ dari $!C, \Gamma \Sequent \Delta$.
  Dengan menerapkan \olref[seq][adm]{prop:weak-G3c-adm} pada $\pi_2''$,
  kita juga memperoleh !!a{proof}~$\pi_2'''$ dari $!C, \Gamma \Sequent
  \Delta, !B$. !!{proof}~$\pi'$ ialah:
  \[
  \AxiomC{}
  \RightLabel{$\pi_1'$}
  \Deduce$\Gamma \fCenter \Delta, !B, !C$
  \AxiomC{}
  \RightLabel{$\pi_2'''$}
  \Deduce$!C, \Gamma \fCenter \Delta, !B$
  \RightLabel{\CutCS}
  \BinaryInf$\Gamma \fCenter \Delta, !B$
  \AxiomC{}
  \RightLabel{$\pi_2'$}
  \Deduce$!B, \Gamma \fCenter \Delta$
  \RightLabel{\CutCS}
  \BinaryInf$\Gamma \fCenter \Delta$
  \DisplayProof
  \]
  \item $!A \ident (!B \land !C)$. Latihan.
  \item $!A \ident (!B \lif !C)$. !!{proof}s $\pi_1$ dan $\pi_2$
  berakhir dengan $\Gamma \Sequent \Delta, !B \lif !C$ dan $!B \lif !C,
  \Gamma \Sequent \Delta$. Menurut \olref{lem:inv-G3c-cut} (keterbalikan
  untuk \RightR{\lif} yang diterapkan pada~$\pi_1$), terdapat !!a{proof}
  $\pi_1'$ dari $!B, \Gamma \Sequent \Delta, !C$. Menurut
  \olref{lem:inv-G3c-cut} (keterbalikan untuk \LeftR{\lif} yang diterapkan
  pada~$\pi_2$), terdapat !!a{proof} $\pi_2'$ dari $\Gamma \Sequent
  \Delta, !B$ dan !!a{proof} $\pi_2''$ dari $!C, \Gamma \Sequent \Delta$.
  Dengan menerapkan \olref[seq][adm]{prop:weak-G3c-adm} pada $\pi_2''$,
  kita juga memperoleh !!a{proof}~$\pi_2'''$ dari $!C, !B, \Gamma \Sequent
  \Delta$. !!{proof}~$\pi'$ ialah:
  \[
  \AxiomC{}
  \RightLabel{$\pi_2'$}
  \Deduce$\Gamma \fCenter \Delta, !B$
  \AxiomC{}
  \RightLabel{$\pi_1'$}
  \Deduce$!B, \Gamma \fCenter \Delta, !C$
  \AxiomC{}
  \RightLabel{$\pi_2'''$}
  \Deduce$!C, !B, \Gamma \fCenter \Delta$
  \RightLabel{\CutCS}
  \BinaryInf$!B, \Gamma \fCenter \Delta$
  \RightLabel{\CutCS}
  \BinaryInf$\Gamma \fCenter \Delta$
  \DisplayProof
  \]
  \item $!A \ident \lforall[x][!B(x)]$. !!{proof}s $\pi_1$ dan
  $\pi_2$ berakhir dengan $\Gamma \Sequent \Delta,
  \lforall[x][!B(x)]$ dan $\lforall[x][!B(x)], \Gamma \Sequent \Delta$.
  Menurut \olref{lem:inv-G3c-cut}, terdapat !!a{proof}~$\pi_1'$ dari
  $\Gamma \Sequent \Delta, !B(c)$.

  Sekarang perhatikan !!{proof}~$\pi_2$. Bukti itu berakhir dengan
  $\lforall[x][!B(x)], \Gamma \Sequent \Delta$. Dengan bergerak ke atas
  dari sekuen akhir~$\pi_2$, tandai setiap kemunculan
  $\lforall[x][!B(x)]$ di sebelah kiri sekuen yang ``menuju ke'' kemunculan
  $\lforall[x][!B(x)]$ dalam sekuen akhir~$\pi_2$, yakni: (a)~Tandai
  $\lforall[x][!B(x)]$ dalam sekuen akhir~$\pi_2$. (b)~Jika
  $\lforall[x][!B(x)]$ muncul dalam konteks kesimpulan suatu aturan dan
  ditandai, tandai kemunculan yang bersesuaian dalam premis-premis aturan
  itu. (c)~Jika $\lforall[x][!B(x)]$ prinsipal dalam kesimpulan aturan
  \LeftR{\lforall} dan ditandai, tandai kemunculan aktif yang bersesuaian
  dari $\lforall[x][!B(x)]$ dan $!B(t)$ dalam premisnya.

  Sekarang perhatikan semua kemunculan $\lforall[x][!B(x)]$ yang ditandai
  tersebut. Tidak satu pun aktif dalam aturan selain \LeftR{\lforall}
  (hanya !!{formula}s aktif dari \LeftR{\lforall} yang ditandai). Hapus
  semua kemunculan $\lforall[x][!B(x)]$ yang ditandai dalam $\pi_2$. Jika
  hasilnya merupakan !!{proof} yang benar, kemunculan-kemunculan yang
  $\lforall[x][!B(x)]$ yang ditandai tadi tidak pernah aktif; semuanya
  berasal langsung dari aksioma.
  !!{proof} itu berakhir dengan~$\Gamma \Sequent \Delta$, sehingga kita
  dapat mengganti $\pi$ dengannya. Kita telah menghapus suatu potongan
  berperingkat maksimal.

  Jika tidak, hasil kita bukan bukti yang benar karena kita telah menghapus
  !!{formula}s $\lforall[x][!B(x)]$ yang aktif dalam beberapa inferensi
  \LeftR{\lforall}. Kita telah mengubah inferensi-inferensi itu menjadi
  ``inferensi'' berbentuk
  \[
  \AxiomC{}
  \RightLabel{$\pi_t$}
  \Deduce$!B(t), \Pi \fCenter \Lambda$
  \RightLabel{\LeftR{\lforall}}
  \UnaryInf$\Pi \fCenter \Lambda$
  \DisplayProof
  \]
  Ganti setiap ``inferensi'' teratas semacam itu dengan
  \[
  \AxiomC{}
  \RightLabel{$\Subst{\pi_1'}{t}{c}[\Pi \Sequent \Lambda]$}
  \Deduce$\Gamma, \Pi \fCenter \Delta, \Lambda, !B(t)$
  \AxiomC{}
  \RightLabel{$\pi_t[\Gamma \Sequent \Delta]$}
  \Deduce$!B(t), \Gamma, \Pi \fCenter \Delta, \Lambda$
  \RightLabel{\CutCS}
  \BinaryInf$\Gamma, \Pi \fCenter \Delta, \Lambda$
  \DisplayProof
  \]
  dan tambahkan $\Gamma$ di sebelah kiri serta $\Delta$ di sebelah kanan
  setiap sekuen di bawahnya. Ulangi sampai semua ``inferensi''
  \LeftR{\lforall} dengan $!B(t)$ yang ditandai dalam premis telah diganti
  oleh inferensi \CutCS{} atas suatu~$!B(t)$. Sekuen akhir !!{proof} yang
  dihasilkan (dan kini benar-benar merupakan !!{proof}) berbentuk $\Gamma,
  \dots, \Gamma \Sequent \Delta, \dots, \Delta$. Terapkan dapat diterimanya
  kontraksi untuk mengubahnya menjadi !!a{proof} dari~$\Gamma \Sequent
  \Delta$, lalu ganti $\pi$ dengannya. Kita telah mengganti satu potongan
  atas $\lforall[x][!B(x)]$ dengan (mungkin banyak) potongan atas $!B(t)$,
  tetapi semuanya berperingkat lebih rendah. Setiap potongan yang telah
  muncul dalam $\pi_1$ atau $\pi_2$ tetap ada, tetapi semuanya juga
  berperingkat lebih rendah.
\end{enumerate}
\end{proof}

\begin{prob}
Periksa kasus $!A \ident (!B \land !C)$ dalam bukti
\olref{lem:max-cut-red-G3c}. Artinya, tunjukkan bahwa jika terdapat
!!a{proof}~$\pi$ yang berakhir dengan
\[
\AxiomC{}
\RightLabel{$\pi_1$}
\Deduce$\Gamma \fCenter \Delta, !B \land !C$
\AxiomC{}
\RightLabel{$\pi_2$}
\Deduce$!B \land !C, \Gamma \fCenter \Delta$
\RightLabel{\CutCS}
\BinaryInf$\Gamma \fCenter \Delta$
\DisplayProof
\]
dengan $\pi_1$ dan $\pi_2$ hanya memuat potongan berperingkat
$<\depth{!B \land !C}$, maka terdapat !!a{proof}~$\pi'$ dari $\Gamma
\Sequent \Delta$ yang hanya memuat potongan berperingkat
$<\depth{!B \land !C}$.
\end{prob}

% \begin{prob}
% Periksa kasus $!A \ident \lexists[x][!B(x)]$ dalam bukti
% \olref{lem:max-cut-red-G3c}.
% \end{prob}


\olref{lem:max-cut-red-G3c} menetapkan bahwa kita dapat mengganti inferensi
\CutCS{} berperingkat tertinggi dalam !!a{proof} dengan inferensi-inferensi
\CutCS{} berperingkat lebih rendah. Sekali lagi, kita dapat menggunakan lema
ini untuk menunjukkan bahwa kita dapat mengeliminasi berapa pun banyaknya
inferensi \CutCS{} dari !!a{proof}, dengan menerapkannya secara berturut-turut
pada sub-!!{proof}s teratas yang berakhir dengan \CutCS{} maksimal.

\begin{thm}\ollabel{thm:cut-elim-G3c-largest}
Setiap !!{proof} $\pi$ dalam $\Log{G3c} + \CutCS$ dapat ditransformasikan
menjadi bukti dalam~\Log{G3c}.
\end{thm}

\begin{proof}
  Pertama, kita buktikan klaim berikut. Misalkan $d$ suatu peringkat yang
  tidak dilampaui oleh peringkat potongan mana pun dalam~$\pi$, dan $n$ banyaknya
  potongan berperingkat~$d$. Semua potongan itu dapat dihapus. Buktinya
  berjalan melalui induksi pada~$n$. Jika $n=0$, tidak ada yang perlu
  dibuktikan. Jika $n>0$, pilih potongan teratas berperingkat~$d$, dan
  misalkan $\pi_1$ sub-!!{proof} yang berakhir dengannya. Menurut
  \olref{lem:max-cut-red-G3c}, terdapat !!a{proof}~$\pi_1'$ atas sekuen akhir
  yang sama, dengan semua potongan berperingkat $<d$. Ganti $\pi_1$ dengan
  $\pi_1'$ dalam~$\pi$. Hasilnya ialah !!a{proof} yang memuat $n-1$
  potongan berperingkat~$d$, sehingga hipotesis induksi berlaku.

  Jika $\pi$ tidak memuat potongan, teorema sudah terbukti. Jika
  tidak, sebut peringkat maksimalnya~$d$. Sekarang hasilnya mengikuti
  melalui induksi pada parameter ini. Jika $d=0$, semua
  potongan atomik, dan menurut hasil sebelumnya semuanya dapat dihapus. Jika
  $d>0$, hasil sebelumnya memberikan suatu bukti tempat semua potongan
  berperingkat $d$ telah diganti oleh potongan berperingkat~$<d$. Karena $d$
  adalah peringkat maksimal potongan dalam~$\pi$, kita memperoleh bukti
  dengan peringkat maksimal potongan $<d$, sehingga hipotesis induksi
  berlaku.
\end{proof}

\end{document}
